\section{Introduction}
\label{sec:intro}

Automatic estimation of student attention is a long-standing objective of learning analytics and affective computing~\cite{whitehill2014faces,chango2022datafusionmmla}. Most systems analyse webcam recordings of individual learners, such as DAiSEE~\cite{gupta2016daisee} and EngageNet~\cite{singh2023engagenet}, and combine several modalities (face, landmarks, head pose and, when available, physiology) through attention or gating~\cite{baltrusaitis2019multimodal,arevalo2017gated,tran2026ordinal}. The underlying premise is that each stream contributes information and that learned fusion weights it appropriately.

This premise is rarely tested in isolation. Reported gains are typically accuracy differences on a single split, obtained with different backbones and protocols for fused and single-stream systems. The multimodal learning literature shows that joint models can underperform their best unimodal branch, because modalities learn at different rates and a dominant modality suppresses the others~\cite{wang2020makes,peng2022balanced}. Whether this holds for classroom attention, where labels are noisy, ordinal and imbalanced, has not been measured.

DIPSER~\cite{marquezcarpintero2025dipser} makes such a test possible. It provides real classroom sessions with per-student cameras and smartwatches, and each window is rated by several annotators. Since a controlled diagnostic requires hundreds of trainings, we encode every stream with discrete Tchebichef moments~\cite{mukundan2001tchebichef}. They require no training and keep each downstream model small, so the study remains affordable without a GPU cluster.

This paper is a diagnostic rather than a state-of-the-art claim. It addresses six questions: (Q1) whether any fusion exceeds the best single stream; (Q2) whether a learnable gate differs from mean pooling; (Q3) whether a linear input-to-head link helps; (Q4) which ordinal objective performs best under rater disagreement; (Q5) whether the answers transfer to DAiSEE and EngageNet; and (Q6) whether they hold on the full DIPSER release, and how much of a reported score depends on the evaluation protocol.

The contributions are as follows. (i) We provide a controlled fusion benchmark on DIPSER, with nested epoch selection and released results for all 495 fits. (ii) We use a shared training-free moment front-end for four modalities, with at most 2.1\,M parameters per model. (iii) We report a negative result with its mechanism: no fusion exceeds the best single stream, and the gate collapses to one stream. (iv) We identify two components that help, the linear link and the soft multi-rater target, and prove that the latter reduces to CORAL under a single label. (v) We provide a transfer study on DAiSEE and EngageNet. (vi) We report a study of the full DIPSER release in which, apart from one exploratory search, success criteria were fixed before each run. It shows that the sensor signal is largely shared by the class, that seven attempts to improve on it fail, and that the evaluation split changes QWK four-fold while accuracy barely changes.
